\documentclass[11pt]{article}
\usepackage[a4paper,margin=2.5cm]{geometry}
\usepackage[T1]{fontenc}
\usepackage[utf8]{inputenc}
\usepackage{amsmath,amssymb}
\usepackage{booktabs,array}
\usepackage{xcolor}
\usepackage{listings}
\usepackage[numbers,sort&compress]{natbib}
\usepackage[colorlinks=true,linkcolor=blue!50!black,citecolor=blue!50!black,urlcolor=blue!50!black]{hyperref}
\lstset{basicstyle=\ttfamily\small,columns=fullflexible,frame=single,framerule=0.3pt,
  rulecolor=\color{black!30},xleftmargin=4pt,xrightmargin=4pt}

\title{Polarized light and geometric phase in qg units:\\
Stokes parameters, Mueller matrices and Stokes' theorem}
\author{Vicente Humberto Monteverde\\
\small Universidad del Museo Social Argentino (UMSA), Argentina\\
\small ORCID 0000-0001-8884-4811}
\date{\small Draft of 3 October 2026}

\begin{document}
\sloppy
\maketitle

\begin{abstract}
The qang project describes qubits by measurement-level quantities, the qg values
$qg_P=\langle P\rangle$~\cite{monteverde2026qang}. This short note states two classical results in those units,
implemented and checked in the \texttt{qang} library (\texttt{pip install qang}, version 0.6.0). First, a
polarization state of light is a qubit: the normalized Stokes parameters are exactly the qg values, the
degree of polarization is the length of the qg vector, and a lossless Mueller matrix is the qg gate rule
$qg'_P=qg_{U^\dagger PU}$. Malus's law becomes $I=I_0(1+qg_Z\cos2\theta+qg_X\sin2\theta)/2$ for any input
state. Second, the geometric phase of a closed loop on the Bloch (Poincar\'e) sphere is minus half the enclosed
solid angle, which is Stokes' theorem for the Berry curvature. For a loop at constant $qg_Z$ the phase is
$(qg_Z-1)/2$ of a turn: it is fixed by $qg_Z$ alone. Nothing here is new physics. The contribution is the
translation into qang's units and a set of 26 numerical checks, with the phase computed three independent
ways.
\end{abstract}

\section{Stokes parameters are qg values}

With $|H\rangle=|0\rangle$ and $|V\rangle=|1\rangle$, the Stokes parameters~\cite{stokes1852} of a polarization
state $\rho$ are $S_\mu=I_0\,\mathrm{Tr}(\rho\,\sigma_\mu)$ with $\sigma=(I,Z,X,Y)$. Hence
\begin{equation}
qg_Z=\frac{S_1}{S_0}=\frac{I_H-I_V}{S_0},\qquad
qg_X=\frac{S_2}{S_0}=\frac{I_D-I_A}{S_0},\qquad
qg_Y=\frac{S_3}{S_0}=\frac{I_R-I_L}{S_0},
\end{equation}
with $|D\rangle,|A\rangle=(|H\rangle\pm|V\rangle)/\sqrt2$ and $|R\rangle,|L\rangle=(|H\rangle\pm i|V\rangle)/\sqrt2$.
Texts differ on which circular state is called right-handed. The library fixes $R$ as above, so $S_3>0$ means
$qg_Y>0$, and offers the opposite convention as an option. The Poincar\'e sphere is the Bloch sphere. Partially
polarized light is a mixed state: its degree of polarization $P=\sqrt{S_1^2+S_2^2+S_3^2}/S_0$ is the length of
the qg vector, and its purity is $\mathrm{Tr}\rho^2=(1+P^2)/2$. A depolarizer
$\mathrm{diag}(1,d,d,d)$ is the depolarizing channel.

\paragraph{Mueller matrices.} For a Jones matrix $J$, $M_{\mu\nu}=\tfrac12\mathrm{Tr}(\sigma_\mu J\sigma_\nu
J^\dagger)$. In the order $(I,Z,X,Y)$ this is the Pauli transfer matrix. For a lossless element (wave plate,
rotator), $J$ is unitary, and the action of $M$ on the Stokes vector is the qg gate rule of
\texttt{qang.formulation}. A half-wave plate at angle $\theta$ sends linear polarization at angle $\alpha$ to
$2\theta-\alpha$, a rotation by $4\theta-2\alpha$ in the $(qg_Z,qg_X)$ plane.

\paragraph{Malus's law in qg units.} Behind an ideal linear polarizer at angle $\theta$, the transmitted
intensity is
\begin{equation}
I(\theta)=\frac{I_0}{2}\left(1+qg_Z\cos2\theta+qg_X\sin2\theta\right),
\end{equation}
which reduces to $I_0\cos^2\theta$ for horizontal light. Since $I(\theta)$ depends on the input only through
$qg_Z$ and $qg_X$, two polarizer settings read two qg values and a quarter-wave plate gives the third.

\paragraph{From counts.} \texttt{qg\_from\_counts} turns analyser counts $(n_H,n_V,n_D,n_A,n_R,n_L)$ into the three
qg values and $P$, with intervals from \texttt{qang.statistics.qg\_estimate} (Bayesian, Wilson or delta method).
Single-photon polarimetry is then the same estimation problem as qubit tomography in qg units.

\section{Geometric phase and Stokes' theorem}

A qubit carried around a closed loop $C$ on the Bloch sphere acquires a geometric
phase~\cite{pancharatnam1956,berry1984}
\begin{equation}
\gamma(C)=-\tfrac12\,\Omega(C),
\end{equation}
where $\Omega$ is the solid angle enclosed by $C$. This is Stokes' theorem. The Berry connection
$A=i\langle\psi|\nabla\psi\rangle$ has curvature $F=\nabla\times A$, a monopole field of strength $1/2$ at the
centre of the sphere. Its line integral around $C$ equals its flux through the enclosed cap, $\Omega/2$. For a
loop at constant $qg_Z=\cos\vartheta$ around the $Z$ axis, $\Omega=2\pi(1-\cos\vartheta)=2\pi(1-qg_Z)$, so in the
qg$_\Phi$ phase unit (fractions of a turn)
\begin{equation}
\phi=\frac{\gamma}{2\pi}=\frac{qg_Z-1}{2}\pmod 1 .
\end{equation}
The geometric phase of a cone loop is fixed by $qg_Z$ alone. For polarized light the same phase is Pancharatnam's
phase on the Poincar\'e sphere. For example, the loop $H\to D\to R\to H$ encloses an octant ($\Omega=\pi/2$) and
gives $\gamma=-\pi/4$.

\paragraph{Three independent computations.} \texttt{qang.geometric} computes the phase three ways:
\begin{enumerate}
\item the gauge-invariant overlap product $\gamma=-\arg\prod_k\langle\psi_k|\psi_{k+1}\rangle$ over a discrete
  loop of states~\cite{samuel1988};
\item the signed solid angle of the loop as a fan of spherical triangles~\cite{vanoosterom1983};
\item the numerical flux of the curvature through the cap.
\end{enumerate}
For a discrete loop the overlap phase equals $-\tfrac12$ the solid angle of the geodesic polygon exactly. The
flux and the line integral agree as Stokes' theorem requires.

\section{Checks}

\begin{table}[h]
\centering\small
\begin{tabular}{p{0.62\textwidth}l}
\toprule
check & result\\
\midrule
six basis states $H,V,D,A,R,L$ map to $\pm1$ on one qg axis & exact\\
Stokes $\leftrightarrow$ qg round trip, against \texttt{qang.formulation.qg\_values} (20 random states) & $10^{-12}$\\
lossless Mueller = qg gate rule (4 elements $\times$ 5 states) & $10^{-12}$\\
Malus's law for $H$, $D$, $R$ and an elliptic state, 7 angles & $10^{-12}$\\
half-wave plate rotates linear polarization to $2\theta-\alpha$ & $10^{-12}$\\
degree of polarization and purity after a depolarizer (10 states) & $10^{-12}$\\
overlap phase $=-\tfrac12$ geodesic solid angle (4 cone heights) & $10^{-10}$\\
cone phase $=(qg_Z-1)/2$ turns (4 heights, 4000 points) & $10^{-5}$\\
curvature flux $=\pi(1-qg_Z)=$ line integral (Stokes, 3 caps) & $10^{-5}$\\
gauge invariance; sign flip under reversal & $10^{-12}$\\
Pancharatnam octant $H\to D\to R$: $\Omega=\pi/2$, $\gamma=-\pi/4$ & exact\\
\bottomrule
\end{tabular}
\caption{Numerical checks in \texttt{tests/test\_polarization.py} (12 tests) and \texttt{tests/test\_geometric.py}
(14 tests).}
\end{table}

\section{Scope}

The results are standard optics and quantum mechanics. Why write them in qg units? The qang tools for qubits
apply unchanged to photonic polarization qubits. These include interval estimation from counts, the gate rule
as a table of qg transformations, and the qg$_\Phi$ phase unit. Maxwell's equations themselves are not part of
this note: they describe classical fields, while qang measures qubit quantities.

\section*{Code}

\begin{lstlisting}[language=Python]
from qang import polarization as P, geometric as G
P.stokes_to_qg([1.0, 0.6, 0.0, 0.8])     # {'X': 0.0, 'Y': 0.8, 'Z': 0.6}
P.malus_intensity({"Z": 1.0}, 0.3)        # cos(0.3)**2
G.cone_phase_turns(0.5)                   # 0.75 turn = -pi/2 rad
\end{lstlisting}

\bibliographystyle{unsrtnat}
\bibliography{../refs}

\end{document}
